\documentclass[12pt]{article}

\usepackage[utf8]{inputenc}
\usepackage[T1]{fontenc}
\usepackage[brazilian]{babel}
\usepackage{sbc-template}
\renewcommand{\ttdefault}{lmtt} % Courier não vem no TinyTeX; Latin Modern Mono no lugar
\usepackage{graphicx}
\usepackage{url}
\usepackage{booktabs}
\usepackage{tabularx}
\usepackage{array}
\usepackage{xcolor}
\usepackage{microtype}
\usepackage[hidelinks]{hyperref}

\sloppy

% Marca o que ainda falta escrever: tem de sumir antes da entrega final.
\newcommand{\pendente}[2]{\textcolor{red!70!black}{\textbf{[#1]} #2}}
\newcolumntype{Y}{>{\raggedright\arraybackslash}X}

\title{Quanto Custa Instalar? Um Experimento com\\
  Seis Gerenciadores de Pacotes Node no Windows 11}

\author{Bernardo S. Alvim\inst{1}, Luísa O. Jardim\inst{1}, Pedro A. S. Seabra\inst{1}}

\address{Engenharia de Software\\
  Pontifícia Universidade Católica de Minas Gerais (PUC Minas)\\
  Belo Horizonte -- MG -- Brasil}

\begin{document}

\maketitle

\begin{abstract}
  \pendente{Luísa, T22b}{Abstract, depois dos resultados.}
\end{abstract}

\begin{resumo}
  \pendente{Luísa, T22b}{Resumo, depois dos resultados.}
\end{resumo}

\section{Introdução}
\label{sec:introducao}

\pendente{Luísa, T22b}{Contexto, os três critérios do software (real, em operação e de alto
impacto, com fontes e datas), as perguntas de pesquisa e as contribuições. Base:
\texttt{docs/proposta.md}, §§2--3.}

\section{Método}
\label{sec:metodo}

\subsection{Objetivo e perguntas de pesquisa}

O estudo combina um experimento controlado com a análise de dados de operação. Seguindo o
GQM~\cite{basili1994gqm}, o objetivo é \textbf{analisar} os gerenciadores de pacotes do ecossistema
Node \textbf{com o propósito de} caracterizá-los e compará-los \textbf{com respeito a} tempo de
instalação, consumo de CPU e memória, pegada em disco e validade da instalação, \textbf{do ponto de
vista de} quem desenvolve e de quem mantém pipelines de integração contínua (CI), \textbf{no
contexto de} três projetos Node, quatro cenários de instalação e três máquinas Windows~11, com um
registry npm local isolando a rede.

As perguntas de pesquisa são:

\begin{itemize}
  \item \textbf{RQ1 (Todos $\times$ npm):} cada gerenciador é mais rápido que o npm, o controle, em
    cada cenário?
  \item \textbf{RQ2 (Yarn $\times$ Yarn):} as três variantes do Yarn diferem entre si?
  \item \textbf{RQ3 (Promessa $\times$ realidade):} o ganho medido confirma o que pnpm e Bun
    anunciam?
  \item \textbf{RQ4 (Adoção $\times$ desempenho):} o gerenciador mais rápido é o mais adotado?
\end{itemize}

\subsection{Desenho experimental}

O experimento tem desenho fatorial $6 \times 4 \times 3$ (gerenciador, cenário e projeto) em blocos,
com cada máquina sendo um bloco. São 72 células por máquina e 10 repetições por célula, número
definido por uma análise de poder sobre uma coleta piloto. Cada máquina executa, portanto, 720
instalações medidas e 2 de aquecimento, descartadas na análise. A ordem das células é sorteada a
cada repetição, com semente derivada do nome da máquina. A Tabela~\ref{tab:variaveis} resume as
variáveis.

\begin{table}[ht]
\centering
\caption{Variáveis do experimento}
\label{tab:variaveis}
\small
\begin{tabularx}{\textwidth}{@{}p{3cm}Y@{}}
\toprule
\textbf{Papel} & \textbf{Variáveis} \\
\midrule
Tratamento & npm 11.19.1 (controle), pnpm 12.21.0, Bun 1.4.4, Yarn Classic 1.22.22, Yarn Berry
4.18.1 com \textit{node-modules} e Yarn Berry 4.18.1 com PnP \\
Fatores & cenário (4) e projeto-carga (3) \\
Bloco & máquina (3) \\
Dependentes & tempo de relógio (principal), CPU, pico de RAM, disco do projeto, cache, número de
arquivos e validade do build \\
Controladas & Node 24.22.3, registry local, mesmo \texttt{package.json} nos seis tratamentos,
Defender no padrão, perfil de alto desempenho, Windows Update pausado e máquina na tomada \\
\bottomrule
\end{tabularx}
\end{table}

Os cenários (Tabela~\ref{tab:cenarios}) definem o que existe antes da instalação medida e seguem o
esquema da página de benchmarks do pnpm, o que permite confrontar os resultados com os números
divulgados.

\begin{table}[ht]
\centering
\caption{Cenários de instalação}
\label{tab:cenarios}
\small
\begin{tabularx}{\textwidth}{@{}lcccY@{}}
\toprule
\textbf{Cenário} & \textbf{Lockfile} & \textbf{Cache} & \textbf{\texttt{node\_modules}} &
\textbf{Situação representada} \\
\midrule
\texttt{cold} & não & não & não & primeira instalação do projeto \\
\texttt{ci} & sim & não & não & pipeline de CI sem cache \\
\texttt{warm-cache} & sim & sim & não & CI com cache restaurado \\
\texttt{reinstall} & sim & sim & sim & instalação repetida com tudo no lugar \\
\bottomrule
\end{tabularx}
\end{table}

\subsection{Objetos}

Os projetos-carga (Tabela~\ref{tab:projetos}) têm dependências em versões exatas e perfis distintos.
Dois vêm de templates oficiais, e a API Express foi escrita pelo grupo, pois o Express não oferece
template oficial de servidor. O objeto de pesquisa é o gerenciador em funcionamento; o projeto é
apenas a carga que ele processa.

\begin{table}[ht]
\centering
\caption{Projetos-carga}
\label{tab:projetos}
\small
\begin{tabularx}{\textwidth}{@{}lccY@{}}
\toprule
\textbf{Projeto} & \textbf{Diretas} & \textbf{Transitivas} & \textbf{Perfil} \\
\midrule
\texttt{express-api} & 4 & 82 & muitos pacotes pequenos \\
\texttt{react-vite} & 9 & 88 & poucos pacotes grandes \\
\texttt{next-app} & 7 & 53 & poucos pacotes muito grandes \\
\bottomrule
\end{tabularx}
\end{table}

\subsection{Instrumentação}

Um coletor próprio em Python monta um diretório limpo com o projeto, prepara o cenário, executa e
mede a instalação, calcula o disco, roda o build e grava uma linha por instalação. A retomada após
interrupção não duplica medições. O tempo de relógio é a métrica principal. CPU e RAM vêm de
amostragem da árvore de processos a cada 50~ms e são métricas secundárias, pois subestimam processos
curtos. A pegada em disco soma os artefatos do projeto e o cache, contando cada inode uma vez para
não duplicar arquivos ligados por \textit{hard link}.

Todos os tratamentos instalam de um registry local (Verdaccio 7.1.0), o que retira da medição a
velocidade da internet de cada casa e evita carga repetida sobre o registry público. Cada gerenciador
tem um cache isolado, apagado antes de cada instalação. Antes de coletar, o coletor confere a versão
de cada ferramenta e pergunta a cada gerenciador onde fica o cache e qual registry será usado,
abortando se houver divergência. Essa checagem foi criada depois que o piloto revelou que o pnpm 12
ignorava em silêncio a configuração recebida; as linhas afetadas foram descartadas e recoletadas.

\subsection{Operação}

A coleta segue um protocolo único nas três máquinas, com ambiente portátil de versões fixadas, apps
fechados, Windows Update pausado e Defender no padrão. As especificações de cada máquina são
registradas antes da coleta (Tabela~\ref{tab:maquinas}).

\begin{table}[ht]
\centering
\caption{Máquinas (blocos)}
\label{tab:maquinas}
\small
\begin{tabularx}{\textwidth}{@{}lYcc@{}}
\toprule
\textbf{Bloco} & \textbf{CPU} & \textbf{RAM} & \textbf{Sistema de arquivos} \\
\midrule
\texttt{bernardo} & AMD Ryzen 5 5500 (6 núcleos, 12 threads) & 15,9 GB & NTFS \\
\texttt{luisa} & \pendente{T12b}{após a coleta} & -- & -- \\
\texttt{pedro} & \pendente{T12c}{após a coleta} & -- & -- \\
\bottomrule
\end{tabularx}
\end{table}

Em paralelo, a série histórica de downloads dos gerenciadores é coletada na API pública do npm, com
no máximo uma requisição por segundo. \pendente{Luísa, T15}{Período e pacotes da série.}

\subsection{Análise}

As hipóteses são bicaudais, com $\alpha = 0{,}05$. Como os tempos são assimétricos, a análise é não
paramétrica. Na RQ1, cada gerenciador é comparado ao npm com o teste de Wilcoxon
pareado~\cite{wilcoxon1945} e correção de Holm~\cite{holm1979} sobre as cinco comparações de cada
combinação de projeto e cenário. Na RQ2, as variantes do Yarn são comparadas com o teste de
Friedman~\cite{friedman1937}, seguido de Wilcoxon com Holm. O tamanho do efeito é dado pelo $\delta$
de Cliff~\cite{cliff1993}, com os limiares de \cite{romano2006}, e pelo
\textit{speedup} (mediana do npm dividida pela mediana do tratamento), com intervalo de confiança de
95\% por \textit{bootstrap}. Só entram nos testes as instalações válidas; a taxa de falhas é
reportada à parte.

Os pares são formados por máquina e repetição. Como repetições na mesma máquina não são independentes
\cite{hurlbert1984}, os testes são interpretados como comparações \textbf{nas três máquinas
estudadas}, e não como generalização para outras máquinas. Por isso, a discussão se apoia no tamanho
do efeito e na consistência entre os blocos.

\section{Resultados}
\label{sec:resultados}

\pendente{Pedro, T22c}{RQ1 a RQ3: descritiva, testes, efeito e figuras.}

\pendente{Bernardo, T14}{RQ2: análise da aba Yarn $\times$ Yarn.}

\pendente{Luísa, T16}{RQ4: adoção $\times$ desempenho.}

\section{Impacto}
\label{sec:impacto}

\pendente{Pedro, T20 e T22c}{Horas e custo de CI, com fontes e o que não se pode afirmar.}

\section{Ameaças à Validade}
\label{sec:ameacas}

\pendente{Bernardo, T17}{Condensar \texttt{docs/ameacas.md} nos quatro tipos~\cite{wohlin2012} depois
dos resultados.}

\section{Reaproveitamento dos Laboratórios}
\label{sec:reaproveitamento}

\pendente{Bernardo, T22a}{O que veio dos Labs 01 a 04 e o que foi feito de novo.}

\section{Conclusão e Trabalhos Futuros}
\label{sec:conclusao}

\pendente{Luísa, T21 e T22b}{Conclusões, limitações e trabalhos futuros (\texttt{docs/ameacas.md},
§5).}

\bibliographystyle{sbc}
\bibliography{referencias}

\end{document}
